% AUTO-GENERATED by scripts/generate_paper_results.py from
% evaluation/partitioning/results.json. Do not hand-edit.
\begin{table}[t]
\centering
\small
\begin{tabular}{lrrrrrr}
\toprule
\textbf{Variant} & \textbf{Violations} & \textbf{Edges retained} & \textbf{Edges cut} & \textbf{Singletons} & \textbf{Largest part.} & \textbf{\# Partitions} \\
\midrule
Naive connected components & 68 & 54 & 0 & 0 & 25 & 32 \\
Contradiction-edge deletion (pre-fix) & 3 & 54 & 0 & 37 & 10 & 69 \\
Constraint-based coloring (production) & 0 & 36 & 18 & 59 & 10 & 85 \\
Weighted-constraint variant (research) & 1 & 46 & 8 & 43 & 10 & 75 \\
\bottomrule
\end{tabular}
\caption{Partitioning-method Pareto trade-off (consistency vs.\
structural connectivity) on the current SMRITI-Reference graph (audit
round P1-4/P1-15; run 20260911_125340, 121 claims,
68 CONTRADICTS / 54
SUPPORTS-REFINES-EQUIVALENT edges). Reading down the ``Violations''
column against ``Edges cut'' shows the trade-off directly: zero
violations costs 18 of the graph's 54 structural
edges (production), one violation recovers 10 of
those 18 cut edges (the weighted-constraint variant, at
8 cut), and accepting all violations (naive connected
components) cuts none. No method dominates on every column
simultaneously -- this is a genuine Pareto frontier, not one variant
strictly beating another. Purity, fragmentation, and ARI/NMI are not
reported --- all three require a ground-truth reference partition
SMRITI-Reference does not have (\S\ref{sec:limitations}).}
\label{tab:partitioning}
\end{table}
